
\section{Дүгнэлт}
\label{sec:conclusion}

Энэхүү судалгаагаар хөрөнгийн тоо болон сургалтын ажиглалтын
тооны харьцаа ($N/T$) бага нөхцөлд ковариацын матрицын
агшилтын үнэлэгчдийн глобал хамгийн бага вариацтай багц
(GMVP)-ын оновчлол дахь гүйцэтгэлийг харьцуулан шинжиллээ.
Түүврийн ковариацын матриц (Sample), тогтмол вариацын (Var),
диагональ (Diag), тогтмол корреляцын (Corr), нэг хүчин зүйлийн
(SF) зорилтот матрицад суурилсан шугаман агшилтын үнэлэгчид
болон шугаман бус агшилтын үнэлэгчийг (NLS) авч үзэв.
Корреляцын бүтэц ялгаатай, тус бүр 10 хувьцаанаас бүрдсэн
гурван багцад богино борлуулалтыг зөвшөөрөөгүй GMVP
байгуулж, 120 сарын сургалтын цонх болон 52 сарын
түүврээс гадуурх (OOS) хугацаанд үнэлэв.

Үндсэн шинжилгээнд агшилтын үнэлэгчдийн гүйцэтгэл
багцын корреляцын бүтцээс хамааран ялгаатай байв.
Өндөр корреляцтай багцад Var үнэлэгч Sample аргатай
харьцуулахад OOS дисперсийг 5.81\%-иар, дунд
корреляцтай багцад Corr үнэлэгч 2.26\%-иар
бууруулсан. Харин бага корреляцтай багцад Sample
арга хамгийн бага OOS дисперстэй байсан бөгөөд
агшилтын таван үнэлэгчийн харьцангуй дисперсийн
бууралт (RVR) бүгд сөрөг гарсан. Иймээс нэг
үнэлэгч бүх багцад тогтмол давуу гүйцэтгэл
үзүүлээгүй.

Багцын жингийн шинжилгээгээр Corr үнэлэгч
гурван багцад хамгийн бага дундаж turnover,
Var үнэлэгч хамгийн бага дундаж HHI үзүүлсэн.
Гэвч жингийн эргэлт болон төвлөрөл бага байх
нь OOS дисперс заавал бага байхыг илэрхийлээгүй.
Тухайлбал, бага корреляцтай багцад Var
үнэлэгч хамгийн бага HHI-тэй боловч OOS
дисперс нь Sample аргынхаас 22.70\%-иар
өндөр байв.

Сургалтын цонхны уртыг 60, 90 болон 120 сар
болгон өөрчилсөн мэдрэмжийн шинжилгээнд
үнэлэгчдийн гүйцэтгэлийн эрэмбэ бүх багцад
тогтвортой хадгалагдсангүй. Өндөр корреляцтай
багцад хамгийн өндөр RVR үзүүлсэн үнэлэгч
сургалтын цонхноос хамааран өөрчлөгдсөн бол
дунд корреляцтай багцад Corr үнэлэгч,
бага корреляцтай багцад Sample арга бүх
цонхонд хамгийн бага OOS дисперстэй байв.
Иймээс $N/T$ харьцаа буурах тусам агшилтын
аргуудын давуу тал заавал нэмэгдэнэ гэсэн
зүй тогтол ажиглагдсангүй.

Үндсэн гурван багцын сарын квадрат өгөөжийн
ялгаанд хийсэн HAC шалгалтаар Holm-ийн
засварын дараа 5\%-ийн түвшинд статистикийн
ач холбогдолтой ялгаа илрээгүй. Сургалтын
цонхны мэдрэмжийн шинжилгээнд мөн ижил
үр дүн гарсан. Эдгээр шалгалт нь квадрат
өгөөжийн дундаж ялгаанд суурилсан тул
OOS дисперсийн ялгааг шууд шалгасан гэж
үзэхгүй.

Үр дүн багцын бүрэлдэхүүнээс хамаарах
эсэхийг үнэлэхийн тулд корреляцын ангилал
тус бүрээс 50, нийт 150 нэмэлт багцад
шинжилгээг давтав. Өндөр корреляцтай
ангилалд агшилтын бүх үнэлэгчийн дундаж
RVR эерэг байсан бол дунд болон бага
корреляцтай ангиллуудад үр дүн харилцан
адилгүй байв. Ялангуяа бага корреляцтай
ангилалд Var үнэлэгчийн дундаж RVR
$-35.02\%$, харин NLS үнэлэгчийнх
$0.99\%$ байв. Иймээс анхны бага
корреляцтай багцад бүх агшилтын
үнэлэгчийн RVR сөрөг гарсан үр дүн
нэмэлт 50 багцад бүрэн давтагдсангүй.

Нэмэлт 150 багцын дундаж RVR-ийн
тодорхой бус байдлыг circular
moving-block bootstrap аргаар үнэлэхэд
бага корреляцтай ангиллын Var
үнэлэгчээс бусад бүх үнэлэгчийн
засварлаагүй 95\%-ийн итгэх завсар
тэгийг агуулсан. Бага корреляцтай
ангиллын Var үнэлэгчийн итгэх завсар
$[-65.61,-14.67]$ байв. Корреляцын
ангилал хоорондын харьцуулалтад
Var үнэлэгчийн High--Low болон
Medium--Low ялгааны засварлаагүй
bootstrap $p$-утга харгалзан
0.0147, 0.0180 байсан боловч
15 харьцуулалтад Holm-ийн засвар
хэрэглэсний дараа 0.2199, 0.2519
болсон. Иймээс ангилал хоорондын
ялгаа олон харьцуулалтын засварын
дараа статистикийн хувьд
батлагдсангүй.

Тэнцүү жинтэй $1/N$ багцыг нэмэлт
жишиг болгон авч үзэхэд үндсэн гурван
багцад ковариацын зургаан үнэлэгчид
суурилсан GMVP стратегиуд бүгд
$1/N$ багцаас бага OOS хэлбэлзэл
үзүүлсэн. Өндөр корреляцтай багцад
Var үнэлэгчийн $1/N$-тэй харьцуулсан
дисперсийн бууралт 39.61\%,
дунд корреляцтай багцад Corr
үнэлэгчийнх 17.63\%, бага
корреляцтай багцад Sample аргынх
26.06\% байв. Гэвч эрсдэлд
тохируулсан өгөөжийн хувьд
GMVP стратегиуд бүх нөхцөлд
давуу байгаагүй. Тухайлбал,
бага корреляцтай багцад
$1/N$ стратегийн жилийн
Sharpe-ийн харьцаа GMVP
стратегиудынхаас өндөр байв.

GMVP болон $1/N$ багцын OOS
дисперсийн харьцаанд хийсэн
circular moving-block bootstrap
шинжилгээгээр дээрх гурван
харьцуулалтын засварлаагүй
95\%-ийн итгэх завсар
харгалзан $[-7.09,62.44]$,
$[-3.95,34.33]$ болон
$[-7.84,51.16]$ байв.
Итгэх завсрууд тэгийг
агуулж, багц тус бүрийн
зургаан харьцуулалтад
Holm-ийн засвар хэрэглэсний
дараа статистикийн ач
холбогдолтой ялгаа
тогтоогдоогүй. Иймээс
GMVP стратегиуд $1/N$
багцаас бага OOS дисперстэй
байсан нь ажиглагдсан
үр дүн боловч статистикийн
хувьд тогтвортой давуу
байдлыг батлахад хангалтгүй
байна.

Арилжааны зардлын мэдрэмжийн
шинжилгээнд багцын дахин
тэнцвэржүүлэлттэй холбоотой
нэг талын 10 суурь нэгжийн
зардлыг тооцоход стратегиудын
Sharpe-ийн харьцаа бага
хэмжээгээр буурсан боловч
үндсэн харьцуулалтын ерөнхий
дүр зураг өөрчлөгдөөгүй.
Гэсэн хэдий ч хамгийн бага
дисперстэй стратеги нь
хамгийн өндөр хуримтлагдсан
өгөөжтэй байх албагүй.
Иймээс ковариацын үнэлэгчийн
практик үр ашгийг зөвхөн
OOS хэлбэлзлээр бус,
арилжааны зардлын дараах
гүйцэтгэлтэй хамтатган
үнэлэх шаардлагатай.

\paragraph{Практик ач холбогдол.}

Судалгааны үр дүн нь бага $N/T$
харьцаатай нөхцөлд ковариацын
матрицын агшилтыг түүврийн
ковариацын матрицаас үргэлж
давуу гэж үзэх үндэслэлгүйг
харуулж байна. Агшилтын
үнэлэгчдийн гүйцэтгэл багцын
корреляцын бүтэц, хөрөнгийн
бүрэлдэхүүн болон сургалтын
цонхны уртаас хамааран
өөрчлөгдөж байна.

Нөгөө талаас GMVP стратегиуд
тэнцүү жинтэй багцаас бага
OOS хэлбэлзэл үзүүлсэн нь
ковариацын матрицад суурилсан
оновчлол эрсдэлийг удирдахад
ашигтай байж болохыг харуулна.
Гэхдээ энэхүү ажиглагдсан
давуу тал статистикийн
хувьд батлагдаагүй бөгөөд
бага хэлбэлзэл нь өндөр
Sharpe-ийн харьцаа эсвэл
өндөр хуримтлагдсан өгөөжийг
заавал илэрхийлэхгүй.

Иймээс хөрөнгө оруулалтын
практикт ковариацын үнэлэгчийг
зөвхөн онолын шинж чанар
эсвэл агшилтын аргын төрлөөр
сонгохын оронд rolling-window
OOS шинжилгээгээр харьцуулж,
дисперс, жингийн эргэлт,
төвлөрөл, эрсдэлд тохируулсан
өгөөж болон арилжааны
зардлыг хамтад нь үнэлэх
нь зүйтэй. Мөн Sample болон
$1/N$ стратегийг харьцуулалтын
жишиг болгон ашиглах нь
оновчлолын болон агшилтын
аргын нэмэлт үр ашгийг
ялган үнэлэхэд чухал юм.

Энэхүү судалгааны үр дүн нь
сонгосон хувьцааны багцууд,
богино борлуулалтыг
зөвшөөрөөгүй оновчлол болон
52 сарын OOS хугацааны
хүрээнд хязгаарлагдана.
Багцуудыг бүрэн хугацааны
өгөгдөлтэй хувьцаануудаас
сонгосон нь сонголтын
хазайлт үүсгэх боломжтой.
Мөн сургалтын цонхны уртыг
өөрчлөхөд $N/T$ харьцаанаас
гадна сургалтын өгөгдлийн
хугацаа өөрчлөгдсөн.
Нэмэлт 150 багц нь бие
даасан санамсаргүй түүвэр
биш бөгөөд bootstrap
шинжилгээнд багцын
бүрэлдэхүүнийг тогтмол
гэж үзсэн. Түүнчлэн
арилжааны зардлын шинжилгээ
нь нэг талын тогтмол
зардлын таамаглалд
суурилсан бөгөөд зах
зээлийн нөлөөлөл болон
бусад бодит арилжааны
хязгаарлалтыг бүрэн
тусгаагүй.

Нэгдсэн дүнгээр бага $N/T$
харьцаатай нөхцөлд
ковариацын матрицын
агшилтын үнэлэгчид GMVP-ийн
OOS дисперсийг бууруулах
боломжтой боловч энэхүү
давуу тал нь бүх багц
болон үнэлэгчид тогтвортой
хадгалагдсангүй. Харин
GMVP стратегиудын $1/N$
багцаас бага хэлбэлзэлтэй
байсан үр дүн нь ковариацын
матрицад суурилсан
оновчлолын боломжит
ач холбогдлыг харуулж
байна. Гэвч статистикийн
шалгалтуудын олон
харьцуулалтын засварын
дараа агшилтын үнэлэгчдийн
Sample аргаас, мөн GMVP
стратегиудын $1/N$ багцаас
давуу гүйцэтгэлийг батлах
хангалттай нотолгоо
илрээгүй. Иймээс ковариацын
матрицын агшилтыг бага
$N/T$ нөхцөлд эрсдэлийг
заавал бууруулах арга бус,
харин багцын бүтэц,
түүврээс гадуурх гүйцэтгэл
болон практик зардалд
тулгуурлан үнэлэх
шаардлагатай хувилбар
гэж дүгнэв.
